\documentclass[../main.tex]{subfiles}

\begin{document}

\ifSubfilesClassLoaded{

    %\tableofcontents
    
    %\setcounter{chapter}{}
}{}

\chapter{ Cartan Criteria }
% 代靖涵 25120222201319

\section{Jordan Decomposition}

\begin{lemma}{}{11-5-2}
Let $x$ be a linear transformation of the complex vector space $V$. Suppose that $x$ has Jordan decomposition $x = d + n$, where $d$ is diagonalisable, $n$ is nilpotent, and $d$ and $n$ commute.

\begin{enumerate}
\item[(a)] There is a polynomial $p(X) \in \mathbb{C}[X]$ such that $p(x) = d$.

\item[(b)] Fix a basis of $V$ in which $d$ is diagonal. Let $\bar{d}$ be the linear map whose matrix with respect to this basis is the complex conjugate of the matrix of $d$. There is a polynomial $q(X) \in \mathbb{C}[X]$ such that $q(x) = \bar{d}$.
\end{enumerate}
\end{lemma}

\begin{exercise}{}{}
    Let $V$ be a vector space, and suppose that $x \in \mathfrak{gl}(V)$ has Jordan decomposition $d+n$. Show that $\operatorname{ad} x : \mathfrak{gl}(V) \to \mathfrak{gl}(V)$ has Jordan decomposition $\operatorname{ad} d + \operatorname{ad} n$.
\end{exercise}
\begin{proof}
    Only need to show that \(  \operatorname{ad}d  \) is semisimple and \(  \operatorname{ad}n  \) is nilpotent, then the proposition follows from the linearity of \(  \operatorname{ad}  \) and the uniqness of Jordan decomposition. 

    \begin{itemize}
        \item Suppose that \(   e_1,\cdots,e_n   \) is the basis in which \(  d  \) is represented diagonal, say \(  \operatorname{diag} \left(  \lambda _1 ,\cdots , \lambda _{n} \right)    \) .  Let \(  E_{ij} \in \mathfrak{gl}\left( V \right)   \) be the maps that  mapping \(  e_{j}  \) to \(  e_{i}  \), that is \(  E_{ij}\left( e_{k} \right)=  \delta _{j}^{k}e_{i}   \) , \(  k=  1,\cdots,n   \) . Then \[
      \begin{aligned}
        \operatorname{ad}\left( d \right)\left( E_{ij} \right)\left( e_{k} \right)&= \left[ d,E_{ij} \right]\left( e_{k} \right)    \\ 
         &= \left( dE_{ij}-E_{ij}d \right)\left( e_{k} \right)\\ 
          &=   d\left( E_{ij}e_{k} \right)- E_{ij}\left( de_{k} \right)\\ 
           &= d\left(  \delta _{j}^{k}e_{i} \right)- E_{ij}\left(  \lambda _{k}e_{k} \right)\\ 
            &=  \delta _{j}^{k} \lambda _{i}e_{i}-  \lambda _{k} \delta _{j}^{k}e_{i}\\ 
             &= \left(  \lambda _{i}- \lambda _{k} \right) \delta _{j}^{k}e_{i}= \left(  \lambda _{i}- \lambda _{j} \right)E_{ij} 
      \end{aligned}  
        \] Which implies that \(  E_{ij}  \) is a eigenvector for \(  \operatorname{ad}\left( d \right)   \) for each \(  i,j=  1,\cdots,n   \). \(  \operatorname{ad}n  \) is diagonal under the basis \(  E_{ij}, i,j=  1,\cdots,n   \)     
        \item Since \(  n  \) is nilpotent, it follows that \(  \operatorname{ad}n  \) is nilpotent. 
        \item  It follows from Jacobi identity that \[
       \begin{aligned}
        \left( \operatorname{ad}d \right)\circ \left( \operatorname{ad}n \right)\left( y \right)&= \left[ d, \left[ n,y \right]  \right]= -\left[ y,\left[ d,n \right]  \right]-\left[ n,\left[ y,d \right]  \right]  
       \end{aligned}     
        \]Since \(  d,n  \) commutes,  \(  \left[ d,n \right]= 0   \).  Thus \(  \operatorname{ad}n  \) and \(  \operatorname{ad}d  \) commute.  
    \end{itemize}
    
\end{proof}
\section{Testing for Solvability}

\begin{exercise}{}{11-5-4}
    Let \(  V  \) be a complex vector space, \(  L  \) be a Lie subalgebra of \(  \mathfrak{gl}\left( V \right)   \).   Suppose that $L$ is solvable. Use Lie's Theorem to show that there is a basis of $V$ in which every element of $L'$ is represented by a strictly upper triangular matrix. Conclude that $\operatorname{tr} xy = 0$ for all $x \in L$ and $y \in L'$.
\end{exercise}
\begin{remark}
    \(  \mathbb{C}   \) can replaced by \(  \mathbb{F}  \) such that \(  \operatorname{char}\mathbb{F}= 0  \), \(  \bar{\mathbb{F}}= \mathbb{F}  \).    
\end{remark}
\begin{proof}
    From Lie's Theorem ,there is a basis of \(  V  \) in which every element of \(  L  \) is upper triangular  , say  \(  \left\{ e_1,\cdots ,e_{n} \right\}  \). Take \(  z \in L^{\prime}   \), \(  z= \left[ x,y \right], x,y \in L   \). Then \[
    z\left( e_{1} \right)= [x,y]\left( e_1 \right)= x\left( y\left( e_1 \right)  \right)-y\left( x\left( e_1 \right)  \right)= 0    
    \]Since \(  e_1  \) is simultaneous an eigenvector of \( x   \) and \(  y  \). Let \(  L_{k}= \operatorname{span}\left\{ e_1,\cdots ,e_{k} \right\}  \) , then \[
  z\left( e_{k} \right)= x\left( y\left( e_{k} \right)  \right)-y\left( x\left( e_{k} \right)  \right)   
    \]    Suppose , \[
    y\left( e_{k} \right)= \sum _{j= 1}^{k}y_{i}e_{j},\quad x\left( e_{k} \right)= \sum _{j= 1}^{k}x_{i}e_{j}  
    \]then \[
  \begin{aligned}
    x\left( y\left( e_{k} \right)  \right)+ \operatorname{span}\left\{ e_1,\cdots ,e_{k-1} \right\}&= x\left( y_{i}e_{k} \right)+ \operatorname{span}\left\{ e_1,\cdots ,e_{k-1} \right\} \\ 
     &= x_{i}y_{i}e_{k}+ \operatorname{span}\left\{ e_1,\cdots ,e_{k-1} \right\}
  \end{aligned}  
    \] Similarly, \[
    y\left( x\left( e_{k} \right)  \right)+ \operatorname{span}\left\{ e_1,\cdots ,e_{k-1} \right\}= x_{i}y_{i}e_{k}+ \operatorname{span}\left\{ e_1,\cdots ,e_{k-1} \right\} 
    \]Thus \[
    z\left( e_{k} \right)\in \operatorname{span}\left\{ e_1,\cdots ,e_{k-1} \right\} 
    \]It follows that \(  z  \) is  nilpotent. Thus every element of \(  L^{\prime}   \) is represented by a strictly upper triangular matrix. 

    The matrix of \(  x  \) and \(  y  \) under the basis \(  \left\{  e_1,\cdots,e_n  \right\}  \) are upper triangular matrix and strictly upper triangular matrix , repectively  , say \(  A  \), \(  B  \). Then \[
    \mathrm{tr}xy= \operatorname{tr}\left( AB \right)= \sum _{k= 1}^{n}a_{ik}b_{kj}= 0 
    \] Since for each \(  i> k  \), \(  a_{ik}= 0  \), and for each \(  i\le k  \), \(  b_{kj}\le 0  \).       
\end{proof}

\begin{proposition}{}{11-5-3}
Let $V$ be a complex vector space and let $L$ be a Lie subalgebra of $\mathfrak{gl}(V)$. If $\mathrm{tr}\,xy = 0$ for all $x, y \in L$, then $L$ is solvable.
\end{proposition}
\begin{proof}
    We shall show that every element of \(  L^{\prime}   \) is nilpotent, then it follows from Engel's Theorem that \(  L^{\prime}   \) is nilpotent, thus \(  L  \) is solvable.
    
    Take any element \(  x \in L^{\prime}   \). Suppose that \(  x  \) has a Jordan decomposition \(  x= d+ n  \). We choose a basis in which \(  d  \) is diagonal \(  d= \operatorname{diag} \left(  \lambda _1 ,\cdots , \lambda _{n} \right)    \) and \(  n  \) is strictly upper triangular. We only need to show that \[
    \sum _{i= 1}^{n}  \lambda _{i}  \bar{\lambda}_{i}= 0
    \]    \[
   \begin{aligned}
    \operatorname{tr}\left( \bar{d}x \right) &= \operatorname{tr}\left( \bar{d} d+  \bar{d}n \right)= \operatorname{tr}\left( \bar{d}d+ \bar{d}n \right)\\ 
     &= \operatorname{tr}\left( \bar{d}d \right)+ 0\\ 
      &= \sum _{i= 1}^{n} \lambda _{i}  \bar{\lambda}_{i}    
   \end{aligned}
    \] Since \(  x\in L^{\prime}   \) ,we write \(  x= \left[ y,z \right]   \), \(  y,z\in L^{\prime}   \). \[
   \begin{aligned}
   \operatorname{tr} \left( \bar{d}x \right) &= \operatorname{tr}\left( \bar{d}\left[ y,z \right]  \right)\\ 
    &= \operatorname{tr}\left( \bar{d}yz- \bar{d}zy \right)= \operatorname{tr}\left( \bar{d}yz- y \bar{d}z\right)   \\ 
     &= \operatorname{tr}\left( \left[ \bar{d},y \right]z   \right) 
   \end{aligned}
    \]  
    
    \begin{note}
        这里我们折腾 \(  \operatorname{ad} \bar{d}  \)而非 \(  \bar{d}  \)的意义在于, 多项式对不变性的保持需要我们对结合代数封闭, 而李代数不一定保持结合代数的封闭性.  
    \end{note}
    Thus to show \(  d= 0  \), we only need to show that \(  \left[ \bar{d},y \right]\in L   \), that is \(  \operatorname{ad}_{\bar{d}}\left( y \right)\in L   \).   \(  \operatorname{ad}x  \) has Jordan decomposition, \(  \operatorname{ad}x= \operatorname{ad}d+ \operatorname{ad} n \). Then from Lemma \ref{lem:11-5-2}  , there exists  \(  q\left( X \right)\in \mathbb{C}\left[ X \right]    \) such that  \[
    q\left( \operatorname{ad}x \right) = \operatorname{ad} \bar{d}
    \]  Since \(  L  \) is \(  \operatorname{ad}x  \)-   invariant, so dose \(  q\left( \operatorname{ad}x \right)   \) and \(  \operatorname{ad}\bar{d}  \), which completes the proof.    
\end{proof}

\begin{theorem}{}{11-5-4}
Let $L$ be a complex Lie algebra. Then $L$ is solvable if and only if $\operatorname{tr}(\operatorname{ad} x \circ \operatorname{ad} y) = 0$ for all $x \in L$ and all $y \in L'$.
\end{theorem}

\begin{note}
    由于\(  L  \)的 sovability等价于任意阶导出子代数\(  L^{\left( k \right) }  \) 的sovability, 因此命题中的 \[
    x \in L,\quad y \in L^{\prime} 
    \]对于任意固定的 \(  k\ge 1,l\ge 0  \) 可以替换为  \[
    x\in L^{\left( k \right) },\quad y \in L^{\left( l \right) },\quad 
    \]
\end{note}

\begin{proof}
    Suppose \(  L  \) is sovable. Then \(  \operatorname{ad}\left( L \right) \subseteq \mathfrak{gl}\left( L \right)    \) is sovable.  Then it follows from exercise \ref{ex:11-5-4} that  \(  \operatorname{tr}\left( \operatorname{ad}x\circ \operatorname{ad}y \right)= 0   \) for all \(   x \in L, y \in L^{\prime}   \).
    
    Conversly, it follows from proposition \ref{pps:11-5-3} that \(  \operatorname{ad}\left( L^{\prime}  \right)   \) is sovable, then so dose \(  L^{\prime}   \) and   \(  L  \). 
\end{proof}

\section{Killing Form}
\begin{definition}{}{}
    Let $L$ be a complex Lie algebra. The Killing form on $L$ is the symmetric bilinear form defined by
$$\kappa(x, y) := \operatorname{tr}(\operatorname{ad} x \circ \operatorname{ad} y) ,\quad \text{for } x, y \in L$$
\end{definition}
\begin{proposition}{}{}
    \[
\kappa([x, y], z) = \kappa(x, [y, z]).
\]

\end{proposition}
\begin{theorem}{Cartan's First Criterion}{}
The complex Lie algebra $L$ is solvable if and only if $\kappa(x,y) = 0$ for all $x \in L$ and $y \in L'$.
\end{theorem}

\begin{proof}
    Just a restatement of theorem \ref{thm:11-5-4}    . 
\end{proof}

\begin{lemma}{}{}
    Suppose that $L$ is a Lie algebra and $I$ is an ideal of $L$. We write $\kappa$ for the Killing form on $L$ and $\kappa_I$ for the Killing form on $I$, considered as a Lie algebra in its own right.
If $x, y \in I$, then $\kappa_I(x,y) = \kappa(x,y)$.
\end{lemma}{}{}

\begin{proof}
   We expand a basis of \(  I  \) to a basis of \(  L  \). Since \[
   \operatorname{ad}x\left( L\right)\subseteq \left[ x,L\right]\subseteq I  
   \] Then \(  \operatorname{ad}x  \) is represented of the form  \[
   \begin{pmatrix} 
       A_{x}&B_{x}\\ 
        O&O 
   \end{pmatrix} 
   \]  \(  \operatorname{ad}y  \) is represented of the form \[
   \begin{pmatrix} 
       A_{y}&B_{y}\\ 
        O& O
   \end{pmatrix} 
   \] We have \[
   \begin{aligned}
    \kappa \left( x,y \right)&=  \operatorname{tr} \left( \begin{pmatrix} 
        A_{x}&B_{x}\\ 
         O&O 
    \end{pmatrix}\begin{pmatrix} 
        A_{y}&B_{y}\\ 
         O&O 
    \end{pmatrix}   \right)   \\ 
     &= \operatorname{tr}\left( \begin{pmatrix} 
         A_{x}A_{y}&A_{x}B_{Y}\\ 
          O&O 
     \end{pmatrix}  \right) \\ 
      &= \operatorname{tr}\left( A_{x}A_{y} \right)\\ 
       &=  \kappa _{I}\left( x,y \right)  
   \end{aligned}
   \]

\end{proof}

\subsection{Testing for Semisimplicity}

\begin{definition}{perpendicular space}{}
    \begin{itemize}
        \item Let $\beta$ be a symmetric bilinear form on a finite-dimensional complex vector space $V$. If $S$ is a subset of $V$, we define the \dfntxt{perpendicular space} to $S$ by
$$S^{\perp} = \{x \in V : \beta(x, s) = 0 \text{ for all } s \in S\}.$$
\item  We say that $\beta$ is \dfntxt{non-degenerate} if $V^{\perp} = 0$; that is, there is no non-zero vector $v \in V$ such that $\beta(v, x) = 0$ for all $x \in V$.

    \end{itemize}
    
\end{definition}

\begin{proposition}{}{}
    If $\beta$ is non-degenerate and $W$ is a vector subspace of $V$, then

\[
\dim W + \dim W^{\perp} = \dim V.
\]
\end{proposition}
\begin{proof}
     Define a linear map \(   \varphi :V\to V^{*}  \), \[
      \varphi \left( v \right)\mapsto  \left( u\mapsto \beta \left( u,v \right)  \right)  
     \] Then \[
     \operatorname{ker} \varphi = \left\{ v\in V: \beta \left( u,v \right)= 0 ,\forall u \in V  \right\}= V^{\perp}= 0
     \]Since \(  V  \) is finite-dimensional , we have \(  \operatorname{dim}V= \operatorname{dim}V^{*}  \). Then \(   \varphi   \) is a isomorphism.   Define \(  W^{0}\subseteq V^{*}  \): \[
     W^{0}= \left\{  \lambda \in V^{*}:  \lambda \left( w \right)= 0, \forall w \in W  \right\}
     \]  We have \[
      \varphi \left( W^{\perp} \right)=  W^{0} ,\quad \operatorname{dim}W^{\perp}= \operatorname{dim}W^{0}
     \] Let \(  i: W\hookrightarrow V  \) the inclusion map , for each \(   \lambda ^{*}\in V^{*}  \),  we define a map \(  i ^{*}: V ^{*}\to W^{*}   \) \[
     i ^{*} \left(  \lambda  \right)\left( w \right)=   \lambda \left( i\left( w \right)  \right) 
     \]   Then \(  \operatorname{Im} i ^{*}= W^{*}  \), \(  \operatorname{ker} i ^{*}= W^{0}  \), we have \[
     \operatorname{dim}W^{0}+ \operatorname{dim}W^{*}= \operatorname{dim}V^{*}
     \]  Then it follows that \[
     \operatorname{dim}W^{\top}+ \operatorname{dim}W^{*}= \operatorname{dim}V^{*}\implies  \operatorname{dim}W^{\perp}+ \operatorname{dim}W = \operatorname{dim}V
     \]Since \(  W,V  \) have finite dimension. 
\end{proof}

\begin{exercise}{}{}
    Suppose \(  I  \) is an ideal of \(  L  \). Shows that \(  I^{\perp} \) is an ideal of \(  L  \).   
\end{exercise}

\begin{proof}
      \[
      I^{\perp}= \left\{ x \in L:   \kappa \left( i,x  \right)= 0, \forall i \in I  \right\}
      \]Take \(  x \in I^{\perp}  \),   \(  y \in L  \),  note that \[
      \left[ y,x \right]\in I^{\perp}\iff  \kappa \left( \left[ x,y \right], i  \right), \quad \forall i \in I  
      \] Since we have \[
      \kappa \left( x,\left[ y, i \right]  \right)= 0 , \forall y \in L, i \in I 
      \] Then  \[
       \kappa \left( [x,y],i \right)= \operatorname{tr} \left( \left[ \operatorname{ad}_{x}, \operatorname{ad}_{y} \right]  \circ \operatorname{ad}_{i}\right)= \operatorname{tr}\left( \operatorname{ad}_{x}\circ \left[ \operatorname{ad}_{y}\circ \operatorname{ad}_{i} \right]  \right)=  \kappa \left( x,\left[ y,i \right]  \right)= 0    
      \] For all \(  i \in I  \). That is \(  \left[ x,y \right]\in I^{\perp}   \). \(  \left[ L, I^{\perp} \right]\subseteq I^{\perp}   \).   
\end{proof}

\begin{corollary}{}{}
    If  \(  L  \)  is semisimple, then \(   \kappa   \) is non-degenerate. 
\end{corollary}

\begin{proof}
    Other wise , \(  L^{\perp}  \) is an nonzero ideal of \(  L  \). Since \(  \left( L^{\perp} \right)^{\prime} \subseteq  L  \)   , we have \[
     \kappa \left( L^{\perp}, \left( L^{\perp} \right)^{\prime}   \right)\subseteq  \kappa \left( L^{\perp}, L \right)= 0  
    \] Then it follows from the Cartan First Criteria that \(  L^{\top}  \) is solvable. Then \(  L^{\perp}  \) is a nonzero solvable   ideal of \(  L  \), which is a contradiction to the Semisimplicity of \(  L  \).  
\end{proof}

\begin{theorem}{Cartan's Second Criterion}{}
The complex Lie algebra $L$ is semisimple if and only if the Killing form $\kappa$ of $L$ is non-degenerate.
\end{theorem}

\section{Derivations of Semisimple Lie Algebras}

\begin{proposition}{}{}
    If \(  L  \) is a finite-dimensional complex semisimple Lie algebra, then \(  \operatorname{ad}L= \operatorname{Der}L  \).  
\end{proposition}
\begin{proof}
 For each \(  x \in L  \), \[
 \begin{aligned}
 \operatorname{ad}x\left[ y,z \right]&= \left[ x,\left[ y,z \right]  \right]\\ 
  &=-[z,[x,y]]-[y,[z,x]]\\ 
   &= [\operatorname{ad}x\left( y \right),z ]+ \left[ y, \operatorname{ad}x\left( z \right)  \right]    
 \end{aligned}
 \] \(  \operatorname{ad}x  \) is a derivation. \(  \operatorname{ad}  \) is a Lie algebra homomorphism from \(  L  \) to \(  \operatorname{Der}L  \).
 For each \(   \delta \in \operatorname{Der}L  \), 
 \[
 \begin{aligned}
 [ \delta , \operatorname{ad}x]y&=  \delta \left[ x,y \right]-\operatorname{ad}x\left(  \delta y \right)\\ 
  &=    [ \delta x,y]+ [x, \delta y]-[x,  \delta y]\\ 
   &= \operatorname{ad}\left(  \delta x \right)\left( y \right)  
 \end{aligned}
 \]Thus \(  [ \delta ,\operatorname{ad}x]= \operatorname{ad}\left(  \delta x \right)   \), the image of \(  \operatorname{ad}  \) is an ideal of \(  \operatorname{Der}L  \).   

 Now we bring in our assumption that \(  L  \) is complex and semisimple. First , note that \(  \operatorname{ad}: L\to \operatorname{Der}L  \) is one-to-one , as \(  \operatorname{ker}\operatorname{ad}= Z\left( L \right)= 0   \), so the lie algebra \(  M\coloneqq \operatorname{ad}L  \) is isomorphic to \(  L  \), therefore it is semisimple as well. We have \[
 \operatorname{dim}M+ \operatorname{dim}M^{\perp}= \operatorname{dim}\operatorname{Der}L
 \]  so our aim is to show that \(  \operatorname{dim}M^{\perp}= 0  \). As \(  M  \) is an  ideal of \(  \operatorname{Der}L  \), the  Killing form \(   \kappa _{M}  \) ofr \(  M  \) is the restriction of the Killing form on \(  \operatorname{Der}L  \). By Cartan's Second Criterion, \(   \kappa _{M}  \) is non-degenerate, so \(  M^{\perp}\cap M= 0  \), and hence \(  \left[ M^{\top},M \right]= 0   \). Thus if \(   \delta \in M^{\perp}  \), and \(  \operatorname{ad}x\in M  \), then \(  \left[  \delta , \operatorname{ad}x \right]= 0   \) for \(  M, M^{\perp}  \) are ideals .     But we saw above that \[
 \left[  \delta , \operatorname{ad}x \right] = \operatorname{ad}\left(  \delta x \right) 
 \]so, for all \(  x \in L  \), we have \(   \delta \left( x \right)= 0   \); in other words, \(   \delta = 0  \).   
\end{proof}

\section{Abstract Jordan Decomposition }

\begin{proposition}{}{}
Let $L$ be a complex Lie algebra. Suppose that $\delta$ is a derivation of $L$ with Jordan decomposition $\delta = \sigma + \nu$, where $\sigma$ is diagonalisable and $\nu$ is nilpotent. Then $\sigma$ and $\nu$ are also derivations of $L$.
\end{proposition}

\begin{proof}
    For \(   \lambda \in \mathbb{C}   \), let \[
    L_{ \lambda }= \left\{ x \in L:\left(  \delta - \lambda  \right)  \right\}
    \] \(  L  \) decomposes as a direct sum of generalised eigenspaces \[
    L= \oplus _{ \lambda }L_{ \lambda }
    \] As \(   \sigma   \) acts diagonalisably, then \(   \lambda   \)-eigenspace of \(   \sigma   \) is   \(  L_{ \lambda }  \)  . Take \(  x \in L_{ \lambda }  \), \(  y\in L_{\mu }  \), then \[
   \begin{aligned}
    \left(  \delta - \left(  \lambda + \mu  \right)   1_{L} \right)[x,y]&= [ \delta x,y]+ [x, \delta y]-  [\lambda x,y]-\left[ x,\mu y \right]  \\ 
     &= [\left(  \delta -  \lambda   1_{L} \right)x,y ]+ [x, \left(  \delta -\mu  1_{L}y \right) ]
   \end{aligned}
    \]\(  \left(  \delta - \lambda  1_{L} \right)^{m+ n}[x,y]   \) is the linear combination of the form \(  [\left(  \delta - \lambda  1_{L} \right)^{i}x, \left(  \delta -\mu  1_{L}y \right)^{j}y ]   \)  , which is zero. Thus \(  \left[ x,y \right]\in L_{ \lambda + \mu }   \) , so \[
     \sigma \left( \left[ x,y \right]  \right)= \left(  \lambda + \mu  \right)\left[ x,y \right]   
    \] which is the same as \[
    \left[  \sigma \left( x \right),y  \right] + \left[ x, \sigma \left( y \right)  \right]= \left[  \lambda x,y \right]+ \left[ x, \mu y \right]   
    \] Thus \(   \sigma   \) is a derivation, and so is  \(   \delta = \sigma = \nu   \).  
\end{proof}

\begin{theorem}{}{}
    Let $L$ be a complex semisimple Lie algebra. Each $x \in L$ can be written uniquely as $x = d+n$, where $d,n \in L$ are such that $\operatorname{ad} d$ is diagonalisable, $\operatorname{ad} n$ is nilpotent, and $[d, n] = 0$. Furthermore, if $y \in L$ commutes with $x$, then $[d, y] = 0$ and $[n, y] = 0$.
\end{theorem}

\begin{proof}
    Let \(  \operatorname{ad}x=  \sigma + \nu   \), where \(   \sigma \in \mathfrak{gl}\left( L \right)   \) is diagonalisable, \(  \nu \in \mathfrak{gl}\left( L \right)   \) is niplpotent, and \(  \left[  \sigma ,\nu  \right]= 0   \). We know that \(   \sigma   \) and \(  \nu   \) are derivations of the semisimple Lie algebra \(  L  \).  And we know that \(  \operatorname{ad}L= \operatorname{Der}L  \), so there exsists   \(  d, n \in L  \) such that \(  \operatorname{ad}d=  \sigma   \) and \(  \operatorname{ad}n = \nu   \). As \(  \operatorname{ad}  \) is injective and \[
    \operatorname{ad}x=  \sigma + \nu = \operatorname{ad}\left( d+ n \right) 
    \] we get that \(  x= d+ n  \).Moreover, \(  \operatorname{ad}\left[ d,n \right]= \left[ \operatorname{ad}d,\operatorname{ad}n \right]= 0    \). The uniqueness of \(  d  \) and \(  n  \) follows from the uniqueness of the Jordan decomposition of \(  \operatorname{ad}x  \). Suppose that \(  y\in L  \) and that \(  \left( \operatorname{ad}x \right)y= 0   \). Since \(   \sigma   \) and \(  \nu   \)              may be expressed as polynomials in \(  \operatorname{ad}x  \). Let \[
    \nu = c_0 1_{L}+ c_1 \operatorname{ad}x+ \cdots + c_{r}\left( \operatorname{ad}x \right)^{r} 
    \] Applying \(  \nu   \) to \(  y  \) we see that \(  \nu \left( y \right)= c_0y   \). But \(  \nu   \) is nilpotent and \(  \nu \left( x \right)= c_0x   \), so \(  c_0= x  \). Thus \(  \nu \left( y \right)= 0   \), and so \(   \sigma \left( y \right)= \left( \operatorname{ad}x-\nu  \right)y= 0    \) also.             
\end{proof}

\begin{definition}{}{}
    Let $L$ be a semisimple Lie algebra and let $\theta: L \to \operatorname{gl}(V)$ be a representation of $L$. Suppose that $x \in L$ has abstract Jordan decomposition $x = d+n$. Then the Jordan decomposition of $\theta(x) \in \operatorname{gl}(V)$ is $\theta(x) = \theta(d) + \theta(n)$.
\end{definition}
\end{document}